% ID: SEL-D0060-OB-E02
% ESTADO: validada. S060.
\textbf{a)} La parábola de $f'$ abre hacia arriba porque su vértice está debajo de sus dos raíces. Por tanto, $f'>0$ en $(-\infty,-1)$ y $(3,\infty)$, donde $f$ crece, y $f'<0$ en $(-1,3)$, donde decrece. Los cambios de signo dan un máximo relativo en $x=-1$ y un mínimo relativo en $x=3$.

\textbf{b)} $g(0)=-2$ y $g'(0)=-6$, pues $g'(x)=-6e^{3x}$. La tangente por $(0,-2)$ con pendiente $-6$ es $\boxed{y=-6x-2}$.
